\SourcePage{113}{116}
\section{Eigenvalues of angular momentum}

In finding the eigenvalues of a particle's angular-momentum projection
onto a chosen direction, one may orient the polar axis along that direction
and employ the spherical-coordinate form of the operator. Using (26.14), the equation
$\hat l_z\psi=l_z\psi$ takes the form
\begin{equation}
  -i\frac{\partial\psi}{\partial\varphi}=l_z\psi.
  \tag{27.1}
\end{equation}
\SourcePage{114}{117}
Its solution is
\[
  \psi=f(r,\theta)e^{il_z\varphi},
\]
where $f(r,\theta)$ is an arbitrary function of $r$ and $\theta$. For
$\psi$ to be single-valued, it must be periodic in $\varphi$, with period
$2\pi$. Hence we find\footnote{The customary notation of the eigenvalues of
the angular-momentum projection by the letter $m$---the same letter that also
denotes the particle mass---cannot in practice lead to confusion.}
\begin{equation}
  l_z=m,\qquad m=0,\ \pm1,\ \pm2,\ldots
  \tag{27.2}
\end{equation}

Thus the eigenvalues $l_z$ are the positive and negative integers, including
zero. We denote the factor depending on $\varphi$ that is characteristic of
the eigenfunctions of $\hat l_z$ by
\begin{equation}
  \Phi_m(\varphi)=\frac{1}{\sqrt{2\pi}}e^{im\varphi}.
  \tag{27.3}
\end{equation}
These functions are normalized so that
\begin{equation}
  \int_0^{2\pi}\Phi_m^*(\varphi)\Phi_{m'}(\varphi)\,d\varphi
  =\delta_{mm'}.
  \tag{27.4}
\end{equation}

The eigenvalues of the $z$ component of the total angular momentum of a system
are evidently likewise positive and negative integers:
\begin{equation}
  L_z=M,\qquad M=0,\ \pm1,\ \pm2,\ldots
  \tag{27.5}
\end{equation}
(this follows because $\hat L_z$ is the sum of the mutually commuting
operators $\hat l_z$ of the individual particles).

Since the direction of the $z$ axis has not been distinguished in advance,
the same result clearly holds for $L_x$, $L_y$, and, in general, for the
component of angular momentum along any direction: each can take only integer
values. At first sight this result may seem paradoxical, particularly if it
is applied to two infinitely close directions. One must bear in mind,
however, that the only common eigenfunction of the operators $\hat L_x$,
$\hat L_y$, and $\hat L_z$ corresponds to the simultaneous values
\[
  L_x=L_y=L_z=0;
\]
in this case the angular-momentum vector, and therefore its projection on
every direction, is zero. If at least one
\SourcePage{115}{118}
of the eigenvalues $L_x$, $L_y$, $L_z$ is nonzero, the corresponding
operators have no common eigenfunctions. Put differently, no state exists in
which definite nonzero values are possessed simultaneously by two or three
angular-momentum components referred to distinct directions; one can therefore
only discuss the integer nature of a single component at a time.

Stationary states of a system that differ only in the value of $M$ have the
same energy. This follows already from general considerations, since the
direction of the $z$ axis has not been distinguished beforehand. Thus the
energy levels of a system with conserved nonzero angular momentum are in any
event degenerate\footnote{This circumstance is a special case of the general
theorem stated in \secref{10} on the degeneracy of levels in the presence of at
least two conserved quantities whose operators do not commute. Here those
quantities are the components of angular momentum.}.

We now turn to finding the eigenvalues of the square of angular momentum and
show how they can be obtained solely from the commutation rules (26.8). Let
$\psi_M$ denote the wave functions of stationary states having the same value
of $\mathbf L^2$, belonging to one degenerate energy level, and differing in
the value of $M$\footnote{Here it is understood that there is no additional
degeneracy by which the same energy corresponds to different values of the
square of angular momentum. This is true for the discrete spectrum (apart
from the so-called “accidental degeneracy” in the Coulomb field; see
\secref{36}), and is, generally speaking, not true for energy levels in the
continuous spectrum. Even when there is additional degeneracy, however, the
eigenfunctions can always be chosen to represent states with definite values
of $\mathbf L^2$; from them one can then select states with the same values of
$E$ and $\mathbf L^2$. The mathematical basis for this is that simultaneous
diagonalization of the matrices of mutually commuting operators is always
achievable. In analogous cases below, for brevity, we shall speak as if
there were no additional degeneracy, with the understanding that, in view of
what has just been said, the results obtained do not in fact depend on this
assumption.}.

First, since the two directions of the $z$ axis are physically equivalent,
every possible positive value $M=|M|$ has a corresponding negative value
$M=-|M|$. Let $L$ (a nonnegative integer) denote the largest possible value
of $|M|$ for a given $\mathbf L^2$. The very existence of such an upper bound
follows from the fact that the difference
$\hat{\mathbf L}^{2}-\hat L_z^{2}=\hat L_x^{2}+\hat L_y^{2}$ is the operator
of the necessarily nonnegative physical quantity $L_x^2+L_y^2$,
\SourcePage{116}{119}
and hence its eigenvalues cannot be negative.

Applying the operator $\hat L_z\hat L_\pm$ to an eigenfunction $\psi_M$ of
$\hat L_z$, and using the commutation rules (26.12), we obtain
\begin{equation}
  \hat L_z\hat L_\pm\psi_M=(M\pm1)\hat L_\pm\psi_M.
  \tag{27.6}
\end{equation}
It follows that $\hat L_\pm\psi_M$ is, apart from a normalization constant,
an eigenfunction corresponding to the value $M\pm1$ of $L_z$:
\begin{equation}
  \psi_{M+1}=\operatorname{const}\cdot\hat L_+\psi_M,
  \qquad
  \psi_{M-1}=\operatorname{const}\cdot\hat L_-\psi_M.
  \tag{27.7}
\end{equation}
If $M=L$ is put in the first of these equalities, then we must have identically
\begin{equation}
  \hat L_+\psi_L=0,
  \tag{27.8}
\end{equation}
since, by definition, there are no states with $M>L$. Applying $\hat L_-$ to
this equality and using (26.13), we obtain
\[
  \hat L_-\hat L_+\psi_L
  =\bigl(\hat{\mathbf L}^{2}-\hat L_z^{2}-\hat L_z\bigr)\psi_L=0.
\]
But since the $\psi_M$ are common eigenfunctions of
$\hat{\mathbf L}^{2}$ and $\hat L_z$,
\[
  \hat{\mathbf L}^{2}\psi_L=\mathbf L^2\psi_L,
  \qquad
  \hat L_z^{2}\psi_L=L^2\psi_L,
  \qquad
  \hat L_z\psi_L=L\psi_L,
\]
so the resulting equation gives
\begin{equation}
  \mathbf L^2=L(L+1).
  \tag{27.9}
\end{equation}

Formula (27.9) determines the required eigenvalues of the square of angular
momentum; $L$ takes all nonnegative integer values. For a given value of $L$,
the angular-momentum component $L_z=M$ can take the values
\begin{equation}
  M=L,L-1,\ldots,-L,
  \tag{27.10}
\end{equation}
that is, $2L+1$ distinct values in all. An energy level corresponding to
angular momentum $L$ is therefore $(2L+1)$-fold degenerate; this is usually
called directional degeneracy of angular momentum. A state with zero angular
momentum, $L=0$ (so that all three of its components vanish), is
nondegenerate. We note that the wave function of such a state is spherically
symmetric. This is already clear from the fact that its change under any
infinitesimal rotation, given by $\hat{\mathbf L}\psi$, vanishes in this case.
\SourcePage{117}{120}
For brevity we shall often use the customary phrase “angular momentum $L$”
of the system, meaning angular momentum whose square equals $L(L+1)$; the
$z$ component of angular momentum is usually called simply its “projection.”

We shall denote the angular momentum of one particle by the lowercase letter
$l$; thus for it we write (27.9) as
\begin{equation}
  \mathbf l^2=l(l+1).
  \tag{27.11}
\end{equation}

Let us calculate the matrix elements of $L_x$ and $L_y$ in the
representation in which, together with the energy, $\mathbf L^2$ and $L_z$
are diagonal (\emph{M.~Born, W.~Heisenberg, P.~Jordan}, 1926). We note that,
since $\hat L_x$ and $\hat L_y$ commute with the Hamiltonian, their matrices
are diagonal with respect to energy; that is, every matrix element for a
transition between states of different energies (and different angular
momenta $L$) is zero. It is therefore sufficient to consider matrix elements
for transitions within the group of states with different values of $M$
corresponding to a single degenerate energy level.

From (27.7) one sees that the matrix of $\hat L_+$ has nonzero elements
only for transitions $M-1\to M$, and that of $\hat L_-$ only for
$M\to M-1$. With this in mind, equating the diagonal matrix elements of the
two sides of (26.13) gives\footnote{In the notation for matrix elements, for brevity, we omit all
indices with respect to which they are diagonal (including the index $L$).}
\[
  L(L+1)
  =\langle M|L_+|M-1\rangle\langle M-1|L_-|M\rangle+M^2-M.
\]
Since the operators $\hat L_x$ and $\hat L_y$ are Hermitian,
\[
  \langle M-1|L_-|M\rangle=\langle M|L_+|M-1\rangle^*,
\]
we rewrite this equality as
\[
  \bigl|\langle M|L_+|M-1\rangle\bigr|^2
  =L(L+1)-M(M-1)=(L-M+1)(L+M),
\]
whence\footnote{The choice of sign in this formula is consistent with the
choice of phase factors in the angular-momentum eigenfunctions.}
\begin{equation}
  \langle M|L_+|M-1\rangle
  =\langle M-1|L_-|M\rangle
  =\sqrt{(L+M)(L-M+1)}.
  \tag{27.12}
\end{equation}
For the nonzero matrix elements of $L_x$ and $L_y$ themselves, this gives
\begin{equation}
  \begin{aligned}
    \langle M|L_x|M-1\rangle
      &=\langle M-1|L_x|M\rangle
        =\frac12\sqrt{(L+M)(L-M+1)},\\
    \langle M|L_y|M-1\rangle
      &=-\langle M-1|L_y|M\rangle
        =-\frac{i}{2}\sqrt{(L+M)(L-M+1)}.
  \end{aligned}
  \tag{27.13}
\end{equation}
\SourcePage{118}{121}
We draw attention to the absence of diagonal elements in the matrices of
$L_x$ and $L_y$. Since a diagonal matrix element gives the mean value of a
quantity in the corresponding state, this means that in states with definite
values of $L_z$ the mean values are $\bar L_x=\bar L_y=0$. Thus, if the
projection of angular momentum on some direction in space has a definite
value, the whole vector $\bar{\mathbf L}$ lies along that same direction.
